\section{Generación condicional funcional: Motif Scaffolding e indicaciones biológicas}

\subsection{De la generación incondicional al control funcional}

El verdadero poder del marco EvoDiff radica en la \textbf{generación condicional funcional} (conditional generation): no solo puede generar proteínas desde cero, sino también guiar el proceso de generación de acuerdo con restricciones funcionales específicas. Esto se debe a la flexibilidad natural del marco de difusión discreta.

La idea central es \textbf{Motif Scaffolding} (diseño de un andamiaje para un motivo funcional):

\begin{importantbox}{Concepto central de Motif Scaffolding}
La función clave de muchas proteínas está determinada por una \textbf{subregión local} de la secuencia (llamada motivo o motivo funcional). Por ejemplo, en una proteína de unión a iones de calcio, solo unos cuantos residuos de aminoácidos específicos participan directamente en la unión física con el ion de calcio.

\textbf{El objetivo de Motif Scaffolding}: dado un motivo funcional conocido, diseñar una secuencia de proteína completamente nueva que «sostenga» (scaffold) ese motivo, de modo que pueda cumplir su función en la posición tridimensional correcta.

Esto es similar a la tarea de \textbf{inpainting} (relleno) en el procesamiento de lenguaje natural: dada una parte del contenido de la secuencia, se le pide al modelo que complete el resto.
\end{importantbox}

\subsection{Inpainting como indicación funcional}

La capacidad de Motif Scaffolding de EvoDiff proviene directamente de su objetivo de entrenamiento. Dado que el modelo aprendió a eliminar el ruido de \textbf{todos los patrones posibles de masking}, adquiere de manera natural las siguientes capacidades:

\begin{enumerate}
    \item Recibir una secuencia \textbf{parcialmente conocida}: la parte conocida es el motivo funcional, y el resto de las posiciones son \texttt{[MASK]}
    \item Eliminar el ruido de forma gradual solo en las posiciones enmascaradas, manteniendo intacta la parte del motivo
    \item Generar finalmente una secuencia de proteína que contiene el motivo original pero cuyo resto es completamente nuevo
\end{enumerate}

Amini compara esta capacidad con una \textbf{«indicación biológica»} (biological prompting): el motivo funcional es el «prompt» que se le da al modelo, y este genera, a partir de ese prompt, una proteína completa que satisface la restricción funcional.

\subsection{Estudio de caso: diseño de una proteína de unión a iones de calcio}

Amini presentó en detalle un caso concreto de diseño funcional:

\textbf{Contexto}: existe de forma natural una proteína cuya función es unirse a iones de calcio (Ca$^{2+}$) y transmitir señales celulares. En la secuencia completa de esta proteína, hay una subsecuencia específica (el motivo marcado en verde) que es directamente responsable de la unión física y la interacción con el ion de calcio.

\textbf{Proceso de diseño}:
\begin{enumerate}
    \item Extraer el fragmento de secuencia del motivo de unión a calcio como «prompt»
    \item Colocar el motivo en la posición correspondiente de la secuencia objetivo, y fijar el resto de las posiciones como \texttt{[MASK]}
    \item Ejecutar el flujo de generación condicional de EvoDiff, con el que el modelo rellena de forma gradual las posiciones enmascaradas
    \item Obtener un diseño de proteína que contiene el motivo original de unión a calcio, pero cuya secuencia global es completamente nueva
\end{enumerate}

\textbf{Validación experimental}:
\begin{itemize}
    \item Se sometió la proteína diseñada por EvoDiff (curva azul) y la proteína natural (curva gris) a un experimento de unión a iones de calcio
    \item Los resultados muestran que la proteína diseñada por IA \textbf{exhibe con éxito capacidad de unión a iones de calcio}
    \item Esto significa que EvoDiff no solo puede generar proteínas viables estructuralmente, sino también proteínas \textbf{viables funcionalmente}
\end{itemize}

\begin{warningbox}{La validación funcional aún se encuentra en una etapa temprana}
Amini reconoció con franqueza que, si bien los experimentos actuales de validación funcional son alentadores, siguen siendo una «prueba de concepto» (proof of concept) preliminar. Del éxito de un solo caso a una capacidad sistemática de diseño funcional aún queda una cantidad considerable de trabajo por delante. El ciclo de validación experimental del diseño de proteínas es largo y costoso, y detrás de cada caso exitoso puede haber múltiples intentos fallidos.
\end{warningbox}

\subsection{Resumen de la sección}

El marco de difusión discreta de EvoDiff admite de manera natural la generación condicional funcional (Motif Scaffolding): al usar un motivo funcional conocido como «indicación biológica», guía al modelo para generar una proteína completamente nueva que contenga dicho motivo. El caso de diseño de la proteína de unión a iones de calcio muestra que la proteína generada por EvoDiff no solo es viable estructuralmente, sino que además logra la función deseada. Esto marca un primer paso importante hacia el diseño de proteínas funcionales impulsado por IA.
